% Complete reviewed Unit 1; editorial master: chapters/ch01/student.md
% Markdown SHA-256: 54f75e982d811b9c305d183ff954bab825f8ed1ca8dbd91f2e4ad42de4d2b126
% This TeX is hand-maintained after the initial one-use assembly.
\newcommand{\runningtitle}{Unit 1. Nouns, the article, and the copula}
% Layout inherited from the reviewed previous course; set a running title per unit.
% This is a small hand-maintained TeX layout, not a conversion pipeline.
\RequirePackage{fix-cm}
\documentclass[11pt,twoside]{article}
\usepackage{fontspec}
\setmainfont{Cambria}[
  Renderer=HarfBuzz,
  BoldFont=cambriab.ttf,
  ItalicFont=cambriai.ttf,
  BoldItalicFont=cambriaz.ttf
]
\setmonofont{lmmono10-regular.otf}[BoldFont=lmmonolt10-bold.otf]
\usepackage[paperwidth=165mm,paperheight=241mm,inner=21mm,outer=19mm,
  top=22mm,bottom=20mm,headheight=14pt,headsep=7mm,footskip=10mm]{geometry}
\usepackage{xcolor,array,booktabs,longtable,tabularx}
\usepackage{enumitem,fancyhdr}
\usepackage[unicode,hidelinks]{hyperref}
\definecolor{ink}{HTML}{000000}
\definecolor{muted}{HTML}{4C5558}
\setlength{\parindent}{0pt}
\setlength{\parskip}{3.5pt}
\AtBeginDocument{\fontsize{11}{14.3}\selectfont}
\setlength{\tabcolsep}{4pt}
\renewcommand{\arraystretch}{1.10}
\setlist[itemize]{leftmargin=1.3em,label=\textbullet,itemsep=2pt,topsep=3pt,parsep=0pt}
\setlist[enumerate]{leftmargin=1.8em,itemsep=3pt,topsep=4pt,parsep=0pt}
\widowpenalty=10000
\clubpenalty=10000
\emergencystretch=1em
\raggedbottom
\pagestyle{fancy}
\fancyhf{}
\providecommand{\runningtitle}{Old English}
\fancyhead[LE,RO]{\small\itshape \runningtitle}
\fancyhead[LO,RE]{\small Old English}
\fancyfoot[LE,RO]{\small\thepage}
\renewcommand{\headrulewidth}{0.3pt}
\newcommand{\booktitle}[2]{
  {\fontsize{9}{11}\selectfont\color{muted}#1\par}
  \vspace{3pt}
  {\fontsize{20}{24}\selectfont\color{ink}#2\par}
  \vspace{6pt}
}
\newcommand{\source}[1]{\par{\footnotesize\color{muted}#1\par}}
\newcommand{\readingtext}[1]{
  \par\addvspace{5pt}
  \begingroup\leftskip=1.5em
    \fontsize{12}{16}\selectfont #1\par
  \endgroup\addvspace{5pt}
}
\makeatletter
\renewcommand{\section}{\@startsection{section}{1}{\z@}
  {-2.4ex \@plus -0.6ex \@minus -0.2ex}{0.8ex \@plus 0.2ex}
  {\normalfont\fontsize{12}{15}\selectfont\bfseries\color{ink}}}
\renewcommand{\subsection}{\@startsection{subsection}{2}{\z@}
  {-1.8ex \@plus -0.5ex \@minus -0.2ex}{0.5ex \@plus 0.2ex}
  {\normalfont\fontsize{11}{14}\selectfont\bfseries\color{ink}}}
\makeatother

\hyphenpenalty=10000
\exhyphenpenalty=10000
\emergencystretch=3em
\newcommand{\itempara}[2]{\par\addvspace{3pt}\begingroup\raggedright\noindent\hangindent=1.8em\hangafter=1\makebox[1.8em][l]{#1.}#2\par\endgroup}
\newcommand{\continuation}[1]{\par\begingroup\raggedright\leftskip=1.8em #1\par\endgroup}
\begin{document}
\raggedright
\booktitle{UNIT 1}{Nouns, the article, and the copula}

\section*{1. Spelling and case}

% teaching: CH01-GRAM-ORTHOGRAPHY CH01-GRAM-GENDER CH01-GRAM-CASES
\textbf{Æ æ} is a single vowel letter. \textbf{Þ þ} (thorn) and \textbf{Ð ð} (eth) represent the sounds written \emph{th} in modern English; the choice between the letters does not by itself distinguish the sounds of \emph{thin} and \emph{this}. This unit follows Cook's use of \textbf{ð}. His hooked \textbf{ę} is an editorial form of \emph{e}, retained in such words as \textbf{męte} and \textbf{męnn}. A macron marks a long vowel or diphthong: \textbf{ā, ǣ, ē, ī, ō, ū, ȳ; ēa, ēo, īe}. Preserve the marks printed in each word; do not add them by analogy.[1]

Nouns have masculine, feminine, or neuter gender, independently of the sex or nature of the thing named. Learn the gender with the noun. A noun's declension determines its case endings; its nominative ending alone does not identify either gender or declension.

\par\begingroup
\setlength{\LTleft}{0pt}\setlength{\LTright}{\fill}
\noindent
\begin{tabular}{@{}>{\raggedright\arraybackslash}p{35mm}>{\raggedright\arraybackslash}p{84mm}@{}}
\toprule
\textbf{Case} & \textbf{Uses introduced here} \\
\midrule
Nominative (N) & Subject; a noun identifying the subject after “be.” \\
Genitive (G) & A dependent expressing possession or another relation, often English “of.” \\
Dative (D) & Recipient; the noun after the selected prepositional constructions below. \\
Accusative (A) & Direct object. \\
\bottomrule
\end{tabular}
\par\endgroup

Singular and plural forms may coincide, and one spelling may serve several cases. Its construction and the agreement of accompanying words determine the analysis. The tables give all four cases, even when their forms are identical.[2]

\section*{2. The article}

% teaching: CH01-GRAM-ARTICLE
The definite article means \textbf{the}. It agrees with its own noun in gender, number, and case. The plural forms are common to all three genders.[3]

\par\begingroup
\setlength{\LTleft}{0pt}\setlength{\LTright}{\fill}
\noindent
\begin{tabular}{@{}>{\raggedright\arraybackslash}p{27mm}>{\raggedright\arraybackslash}p{21mm}>{\raggedright\arraybackslash}p{21mm}>{\raggedright\arraybackslash}p{21mm}>{\raggedright\arraybackslash}p{21mm}@{}}
\toprule
\textbf{Case} & \textbf{Masculine singular} & \textbf{Neuter singular} & \textbf{Feminine singular} & \textbf{All genders plural} \\
\midrule
N & se & ðæt & sēo & ðā \\
G & ðæs & ðæs & ðǣre & ðāra \\
D & ðǣm & ðǣm & ðǣre & ðǣm \\
A & ðone & ðæt & ðā & ðā \\
Instrumental: reference & ðȳ, ðon & ðȳ, ðon & as dative & as dative \\
\bottomrule
\end{tabular}
\par\endgroup

Recognize \textbf{ðȳ} and \textbf{ðon} as instrumental singular forms; their uses are not required in composition here. Nouns have no separate instrumental forms in these paradigms. The source also permits \textbf{ðām} for \textbf{ðǣm} and \textbf{ðǣra} for \textbf{ðāra}; the tables and exercises select \textbf{ðǣm} and \textbf{ðāra}.

\section*{3. Strong masculine nouns}

% teaching: CH01-PAR-MASC CH01-PAR-MASC-E
\textbf{Fisc} illustrates the common consonant pattern. In \textbf{męte}, the nominative \textbf{-e} is replaced by \textbf{-es, -as, -a, -um} where those endings occur; it is not part of a stem to which another full ending is added.[4]

\par\begingroup
\setlength{\LTleft}{0pt}\setlength{\LTright}{\fill}
\noindent
\begin{tabular}{@{}>{\raggedright\arraybackslash}p{35mm}>{\raggedright\arraybackslash}p{40mm}>{\raggedright\arraybackslash}p{40mm}@{}}
\toprule
\textbf{Case and number} & \textbf{fisc} & \textbf{męte} \\
\midrule
N singular & fisc & męte \\
G singular & fisces & mętes \\
D singular & fisce & męte \\
A singular & fisc & męte \\
N plural & fiscas & mętas \\
G plural & fisca & męta \\
D plural & fiscum & mętum \\
A plural & fiscas & mętas \\
\bottomrule
\end{tabular}
\par\endgroup

The plural of \textbf{męte} is displayed for recognition of the pattern; this unit uses its singular “food.” \textbf{Heofon, gāst, wæstm} follow \textbf{fisc}. Keep both syllables of \textbf{heofon} before endings. The divine name \textbf{God} has the singular \textbf{God, Godes, Gode, God}; no plural of this proper-name sense is required.

% teaching: CH01-PAR-MASC-AE-A
In \textbf{dæg} and \textbf{hwæl}, singular \textbf{æ} alternates with plural \textbf{a}. The change does not apply to \textbf{wæstm}, whose \textbf{æ} stands before a consonant group.

\par\begingroup
\setlength{\LTleft}{0pt}\setlength{\LTright}{\fill}
\noindent
\begin{tabular}{@{}>{\raggedright\arraybackslash}p{35mm}>{\raggedright\arraybackslash}p{40mm}>{\raggedright\arraybackslash}p{40mm}@{}}
\toprule
\textbf{Case and number} & \textbf{dæg} & \textbf{hwæl} \\
\midrule
N singular & dæg & hwæl \\
G singular & dæges & hwæles \\
D singular & dæge & hwæle \\
A singular & dæg & hwæl \\
N plural & dagas & hwalas \\
G plural & daga & hwala \\
D plural & dagum & hwalum \\
A plural & dagas & hwalas \\
\bottomrule
\end{tabular}
\par\endgroup

% teaching: CH01-PAR-MASC-SYNCOPE
\textbf{Fugol} and \textbf{morgen} lose their unstressed middle vowel before endings in the selected forms. This loss, called \emph{syncope}, is lexical: use the diagnostic forms given for each noun rather than deleting vowels from every word of two syllables. Forms retaining the vowel also occur.[4]

\par\begingroup
\setlength{\LTleft}{0pt}\setlength{\LTright}{\fill}
\noindent
\begin{tabular}{@{}>{\raggedright\arraybackslash}p{35mm}>{\raggedright\arraybackslash}p{40mm}>{\raggedright\arraybackslash}p{40mm}@{}}
\toprule
\textbf{Case and number} & \textbf{fugol} & \textbf{morgen} \\
\midrule
N singular & fugol & morgen \\
G singular & fugles & morgnes \\
D singular & fugle & morgne \\
A singular & fugol & morgen \\
N plural & fuglas & morgnas \\
G plural & fugla & morgna \\
D plural & fuglum & morgnum \\
A plural & fuglas & morgnas \\
\bottomrule
\end{tabular}
\par\endgroup

\section*{4. Strong neuter nouns}

% teaching: CH01-GRAM-WEIGHT CH01-PAR-NEUT-WEIGHT
Neuter nominative and accusative are identical in each number. The genitive and dative endings match those of \textbf{fisc}. In the ordinary monosyllabic pattern, the nominative/accusative plural has \textbf{-u} after a short stem, but no ending after a heavy stem. A short vowel followed by one consonant gives the short pattern here; a long vowel or long diphthong, or a short vowel followed by two consonants, gives the heavy pattern. A doubled consonant counts as two. Thus \textbf{hof} has \textbf{hofu}, whereas \textbf{sǣd, gærs, cynn} have no added N/A plural ending. Vowel length and syllable weight are not the same thing.[5]

\par\begingroup
\setlength{\LTleft}{0pt}\setlength{\LTright}{\fill}
\noindent
\begin{tabular}{@{}>{\raggedright\arraybackslash}p{35mm}>{\raggedright\arraybackslash}p{40mm}>{\raggedright\arraybackslash}p{40mm}@{}}
\toprule
\textbf{Case and number} & \textbf{hof} & \textbf{sǣd} \\
\midrule
N singular & hof & sǣd \\
G singular & hofes & sǣdes \\
D singular & hofe & sǣde \\
A singular & hof & sǣd \\
N plural & hofu & sǣd \\
G plural & hofa & sǣda \\
D plural & hofum & sǣdum \\
A plural & hofu & sǣd \\
\bottomrule
\end{tabular}
\par\endgroup

The heavy pattern also covers \textbf{lēoht, līf, gēar, dēor, anginn, lyft, weorc}. In \textbf{anginn}, apply the rule to the stressed stem \textbf{-ginn}, with \textbf{nn}. Plural forms of mass or abstract nouns do not automatically have the same sense as the singular; the vocabulary below marks those whose plurals are only for recognizing the morphology.

% teaching: CH01-PAR-NEUT-TREO
\textbf{Trēo} has \textbf{w} before every ending and the exceptional N/A plural \textbf{trēowu}, despite its long vowel. Bare singular \textbf{trēow} also occurs.

\par\begingroup
\setlength{\LTleft}{0pt}\setlength{\LTright}{\fill}
\noindent
\begin{tabular}{@{}>{\raggedright\arraybackslash}p{35mm}>{\raggedright\arraybackslash}p{84mm}@{}}
\toprule
\textbf{Case and number} & \textbf{trēo} \\
\midrule
N singular & trēo \\
G singular & trēowes \\
D singular & trēowe \\
A singular & trēo \\
N plural & trēowu \\
G plural & trēowa \\
D plural & trēowum \\
A plural & trēowu \\
\bottomrule
\end{tabular}
\par\endgroup

% teaching: CH01-PAR-NEUT-DISYLLABLE
Do not apply the monosyllabic \textbf{-u/zero} rule mechanically to nouns of two syllables. Learn the following complete patterns. \textbf{Tācen} loses \textbf{e} before endings; \textbf{ǣfen, wæter, nīeten} retain it in the selected forms. The shorter \textbf{wætru} and related forms also occur.[5]

\par\begingroup
\setlength{\LTleft}{0pt}\setlength{\LTright}{\fill}
\noindent
\begin{tabular}{@{}>{\raggedright\arraybackslash}p{27mm}>{\raggedright\arraybackslash}p{21mm}>{\raggedright\arraybackslash}p{21mm}>{\raggedright\arraybackslash}p{21mm}>{\raggedright\arraybackslash}p{21mm}@{}}
\toprule
\textbf{Case and number} & \textbf{ǣfen} & \textbf{tācen} & \textbf{wæter} & \textbf{nīeten} \\
\midrule
N singular & ǣfen & tācen & wæter & nīeten \\
G singular & ǣfenes & tācnes & wæteres & nīetenes \\
D singular & ǣfene & tācne & wætere & nīetene \\
A singular & ǣfen & tācen & wæter & nīeten \\
N plural & ǣfenu & tācnu & wæteru & nīetenu \\
G plural & ǣfena & tācna & wætera & nīetena \\
D plural & ǣfenum & tācnum & wæterum & nīetenum \\
A plural & ǣfenu & tācnu & wæteru & nīetenu \\
\bottomrule
\end{tabular}
\par\endgroup

\section*{5. Strong feminine nouns}

% teaching: CH01-PAR-FEM-LONG CH01-PAR-FEM-ACC-ZERO
The strong feminines required here have no nominative singular ending. The ordinary genitive, dative, and accusative singular end in \textbf{-e}. \textbf{Wyrt} and \textbf{tīd}, however, have an endingless accusative singular in the selected pattern. The selected N/A plural is \textbf{-a}; \textbf{-e} is also found.[6]

\par\begingroup
\setlength{\LTleft}{0pt}\setlength{\LTright}{\fill}
\noindent
\begin{tabular}{@{}>{\raggedright\arraybackslash}p{32mm}>{\raggedright\arraybackslash}p{28mm}>{\raggedright\arraybackslash}p{28mm}>{\raggedright\arraybackslash}p{28mm}@{}}
\toprule
\textbf{Case and number} & \textbf{gesceaft} & \textbf{wyrt} & \textbf{tīd} \\
\midrule
N singular & gesceaft & wyrt & tīd \\
G singular & gesceafte & wyrte & tīde \\
D singular & gesceafte & wyrte & tīde \\
A singular & gesceafte & wyrt & tīd \\
N plural & gesceafta & wyrta & tīda \\
G plural & gesceafta & wyrta & tīda \\
D plural & gesceaftum & wyrtum & tīdum \\
A plural & gesceafta & wyrta & tīda \\
\bottomrule
\end{tabular}
\par\endgroup

% teaching: CH01-PAR-FEM-NES CH01-PAR-FEM-UNG
Nouns ending in \textbf{-nes} have \textbf{-ness-} before endings: retain the \textbf{e} and double the \textbf{s}. The same pattern governs \textbf{drȳgnes, gelīcnes, fæstnes, nēowolnes, brādnes}. Nouns in \textbf{-ung}, and the assigned \textbf{līehting}, take the endings shown for \textbf{gegaderung}; \textbf{frætwung} follows it. Cook also records singular oblique \textbf{-a} in this group, especially in the dative. The present tables select \textbf{-e}.[6]

\par\begingroup
\setlength{\LTleft}{0pt}\setlength{\LTright}{\fill}
\noindent
\begin{tabular}{@{}>{\raggedright\arraybackslash}p{35mm}>{\raggedright\arraybackslash}p{40mm}>{\raggedright\arraybackslash}p{40mm}@{}}
\toprule
\textbf{Case and number} & \textbf{andlīcnes} & \textbf{gegaderung} \\
\midrule
N singular & andlīcnes & gegaderung \\
G singular & andlīcnesse & gegaderunge \\
D singular & andlīcnesse & gegaderunge \\
A singular & andlīcnesse & gegaderunge \\
N plural & andlīcnessa & gegaderunga \\
G plural & andlīcnessa & gegaderunga \\
D plural & andlīcnessum & gegaderungum \\
A plural & andlīcnessa & gegaderunga \\
\bottomrule
\end{tabular}
\par\endgroup

\section*{6. Weak nouns}

% teaching: CH01-PAR-WEAK
Weak masculine \textbf{steorra} and weak feminine \textbf{eorðe} differ in the nominative singular. Both have \textbf{-an} in the other singular cases and the N/A plural, \textbf{-ena} in the genitive plural, and \textbf{-um} in the dative plural. Replace nominative \textbf{-a/-e} with these endings.[7]

\par\begingroup
\setlength{\LTleft}{0pt}\setlength{\LTright}{\fill}
\noindent
\begin{tabular}{@{}>{\raggedright\arraybackslash}p{35mm}>{\raggedright\arraybackslash}p{40mm}>{\raggedright\arraybackslash}p{40mm}@{}}
\toprule
\textbf{Case and number} & \textbf{steorra} & \textbf{eorðe} \\
\midrule
N singular & steorra & eorðe \\
G singular & steorran & eorðan \\
D singular & steorran & eorðan \\
A singular & steorran & eorðan \\
N plural & steorran & eorðan \\
G plural & steorrena & eorðena \\
D plural & steorrum & eorðum \\
A plural & steorran & eorðan \\
\bottomrule
\end{tabular}
\par\endgroup

The plural of \textbf{eorðe} is displayed for recognizing its pattern; the exercises use singular “earth.” An \textbf{-an} form alone cannot identify either case or number.

\section*{7. Individual paradigms}

% teaching: CH01-PAR-EXCEPTION-MANN CH01-PAR-EXCEPTION-NIHT
\textbf{Mann} has \textbf{męnn} in the dative singular and N/A plural. In the selected \textbf{niht} paradigm, every singular case and the N/A plural are \textbf{niht}; the genitive and dative plurals are \textbf{nihta, nihtum}. Genitive/dative singular \textbf{nihte} is a valid alternative, but is not the form selected for recall here.[8]

\par\begingroup
\setlength{\LTleft}{0pt}\setlength{\LTright}{\fill}
\noindent
\begin{tabular}{@{}>{\raggedright\arraybackslash}p{35mm}>{\raggedright\arraybackslash}p{40mm}>{\raggedright\arraybackslash}p{40mm}@{}}
\toprule
\textbf{Case and number} & \textbf{mann} & \textbf{niht} \\
\midrule
N singular & mann & niht \\
G singular & mannes & niht \\
D singular & męnn & niht \\
A singular & mann & niht \\
N plural & męnn & niht \\
G plural & manna & nihta \\
D plural & mannum & nihtum \\
A plural & męnn & niht \\
\bottomrule
\end{tabular}
\par\endgroup

% teaching: CH01-PAR-EXCEPTION-SAE CH01-REF-SAE
\textbf{Sǣ} is masculine in the selected vocabulary. Its nominative plural is \textbf{sǣs}, not a form in \textbf{-as}. The starred genitive plural below is a reconstruction: consult it only as a reference, and do not use it in these exercises.[8]

\par\begingroup
\setlength{\LTleft}{0pt}\setlength{\LTright}{\fill}
\noindent
\begin{tabular}{@{}>{\raggedright\arraybackslash}p{35mm}>{\raggedright\arraybackslash}p{84mm}@{}}
\toprule
\textbf{Case and number} & \textbf{sǣ} \\
\midrule
N singular & sǣ \\
G singular & sǣs \\
D singular & sǣ \\
A singular & sǣ \\
N plural & sǣs \\
G plural & *sǣwa — reconstructed; not tested \\
D plural & sǣum \\
A plural & sǣs \\
\bottomrule
\end{tabular}
\par\endgroup

% teaching: CH01-PAR-EXCEPTION-THIESTRU
\textbf{Ðīestru}, “darkness,” is plural in Old English. Its complete assigned paradigm is plural only.

\par\begingroup
\setlength{\LTleft}{0pt}\setlength{\LTright}{\fill}
\noindent
\begin{tabular}{@{}>{\raggedright\arraybackslash}p{35mm}>{\raggedright\arraybackslash}p{84mm}@{}}
\toprule
\textbf{Case} & \textbf{Plural} \\
\midrule
N & ðīestru \\
G & ðīestra \\
D & ðīestrum \\
A & ðīestru \\
\bottomrule
\end{tabular}
\par\endgroup

\section*{8. The indicative copula}

% teaching: CH01-GRAM-COPULA
These are the complete present \textbf{eom} and preterite \textbf{wæs} series of the indicative “be.” The preterite expresses a past state. Other present forms, negative forms, moods, and nonfinite constructions will be treated separately.[9]

\par\begingroup
\setlength{\LTleft}{0pt}\setlength{\LTright}{\fill}
\noindent
\begin{tabular}{@{}>{\raggedright\arraybackslash}p{35mm}>{\raggedright\arraybackslash}p{40mm}>{\raggedright\arraybackslash}p{40mm}@{}}
\toprule
\textbf{Person and number} & \textbf{Present} & \textbf{Preterite} \\
\midrule
1 singular: I & eom & wæs \\
2 singular: you & eart & wǣre \\
3 singular: he, she, it & is & wæs \\
1 plural: we & sindon & wǣron \\
2 plural: you & sindon & wǣron \\
3 plural: they & sindon & wǣron \\
\bottomrule
\end{tabular}
\par\endgroup

English pronouns label the persons; no Old English pronoun forms are required for this table. Cook also gives \textbf{sind, sint} beside \textbf{sindon}.

\section*{9. Basic constructions}

% teaching: CH01-GRAM-SUBJECT-PREDICATE CH01-GRAM-AGREEMENT
The subject is nominative. A noun that identifies it after “be” is also nominative; English \textbf{a/an} need not receive a separate word in such a predicate. The copula agrees with the subject in person and number.[10]

% example: CH01-EXAMPLE07
\textbf{Se hwæl is dēor.} — The whale is an animal.

% teaching: CH01-GRAM-COORDINATION
\textbf{And} joins words or phrases. When definite noun phrases are joined, each article agrees with its own noun. Gender belongs to that noun; it is not taken from the other noun in the pair.

% teaching: CH01-GRAM-ON
For the location of an entity, use \textbf{on + dative}, “in/on.”

% example: CH01-EXAMPLE01
\textbf{Se fisc is on ðǣm wætere.} — The fish is in the water.

% example: CH01-EXAMPLE02
\textbf{Ðā męnn wǣron on ðǣm hofum.} — The men were in the dwellings.

% teaching: CH01-GRAM-GENITIVE
A dependent genitive expresses possession or a relation such as the source, contents, or kind of something. In the phrase pattern used for Cook's Exercise 3, place the genitive first and omit the article of the \textbf{head noun}, while retaining the article belonging to the dependent. The head remains in the case required by its own role. A proper name such as \textbf{God} takes no article here.[11]

% example: CH01-EXAMPLE03
\textbf{ðæs mannes fugol} — the man's bird

Here \textbf{ðæs mannes} is the dependent genitive; \textbf{fugol} is the head, nominative in the isolated naming phrase. The construction does not mean that every English “the” disappears.

% teaching: CH01-GRAM-NESTED-GENITIVE
The same dependent may stand inside a prepositional phrase. Its genitive does not change the case required for the head.

% example: CH01-EXAMPLE08
\textbf{on ðæs mannes hofe} — in the man's dwelling

\textbf{Hofe} is dative after locative \textbf{on}; \textbf{ðæs mannes} remains genitive. \textbf{Ðæs} agrees with \textbf{mannes}, not with \textbf{hofe}.

% teaching: CH01-GRAM-OBJECT
The direct object is accusative. In the following example, \textbf{geseah} is supplied: “saw,” third-person singular preterite. Its conjugation is not required here.[12]

% example: CH01-EXAMPLE04
\textbf{Se mann geseah ðone fisc.} — The man saw the fish.

\textbf{Mann} and \textbf{fisc} have identical nominative and accusative singular spellings, but their articles distinguish subject \textbf{se} from object \textbf{ðone}.

% teaching: CH01-GRAM-RECIPIENT CH01-FRAME-GIVING
The recipient of a gift is dative; the thing given is accusative. The supplied \textbf{sealde} means “gave,” third-person singular preterite.[12]

% example: CH01-EXAMPLE05
\textbf{God sealde ðǣm fugle męte.} — God gave the bird food.

\textbf{Ðǣm fugle} is the recipient; \textbf{męte} is the direct object. English may express the recipient with “to,” but this Old English construction uses the dative without a preposition.

% teaching: CH01-GRAM-FRAM
Use \textbf{fram + dative} for what something is separated from. The supplied \textbf{gedǣlde} means “separated/divided,” third-person singular preterite.[12]

% example: CH01-EXAMPLE06
\textbf{God gedǣlde ðæt lēoht fram ðǣm ðīestrum.} — God separated the light from the darkness.

The thing separated, \textbf{ðæt lēoht}, is accusative; \textbf{ðǣm ðīestrum} is dative after \textbf{fram}. These are particular constructions. They do not establish one case rule for all spatial prepositions.

\section*{10. Vocabulary}

Learn each noun with its gender and diagnostic forms. \textbf{G sg} identifies the singular stem; \textbf{N pl} identifies the plural pattern. The full tables and rules above supply the other cases. \textbf{S} marks a noun used only in the singular in composition here: its plural morphology is for recognition, with no additional plural sense assigned. \textbf{P} marks plural-only \textbf{ðīestru}. The singular divine name \textbf{God} has no plural assigned. Abbreviations: \textbf{m.}, masculine; \textbf{f.}, feminine; \textbf{n.}, neuter.[13]

% Vocabulary reflow: all six editorial fields retained in four columns.
\par\begingroup
\fontsize{10.5}{13.2}\selectfont
\setlength{\LTleft}{0pt}\setlength{\LTright}{\fill}
\begin{longtable}{@{}>{\raggedright\arraybackslash}p{21mm}>{\raggedright\arraybackslash}p{35mm}>{\raggedright\arraybackslash}p{33mm}>{\raggedright\arraybackslash}p{27mm}@{}}
\toprule
\textbf{Noun; gender} & \textbf{Meaning} & \textbf{Diagnostic forms} & \textbf{Pattern; limit} \\
\midrule
\endfirsthead
\toprule
\textbf{Noun; gender} & \textbf{Meaning} & \textbf{Diagnostic forms} & \textbf{Pattern; limit} \\
\midrule
\endhead
\midrule
\multicolumn{4}{r}{\footnotesize Continued on the following page.}\\
\endfoot
\bottomrule
\endlastfoot
% teaching: CH01-VOC-aefen
ǣfen (n.) & evening & G sg ǣfenes\newline N pl ǣfenu & §4 \\
% teaching: CH01-VOC-morgen
morgen (m.) & morning & G sg morgnes\newline N pl morgnas & fugol/morgen \\
% teaching: CH01-VOC-daeg
dæg (m.) & day & G sg dæges\newline N pl dagas & dæg/hwæl \\
% teaching: CH01-VOC-niht
niht (f.) & night & G sg niht\newline N pl niht & niht \\
% teaching: CH01-VOC-eorthe
eorðe (f.) & earth & G sg eorðan\newline N pl eorðan & steorra/eorðe; S \\
% teaching: CH01-VOC-heofon
heofon (m.) & heaven; sky & G sg heofones\newline N pl heofonas & fisc \\
% teaching: CH01-VOC-fisc
fisc (m.) & fish & G sg fisces\newline N pl fiscas & fisc \\
% teaching: CH01-VOC-fugol
fugol (m.) & bird & G sg fugles\newline N pl fuglas & fugol/morgen \\
% teaching: CH01-VOC-gaers
gærs (n.) & grass & G sg gærses\newline N pl gærs & hof/sǣd; S \\
% teaching: CH01-VOC-saed
sǣd (n.) & seed & G sg sǣdes\newline N pl sǣd & hof/sǣd \\
% teaching: CH01-VOC-treo
trēo (n.) & tree & G sg trēowes\newline N pl trēowu & trēo \\
% teaching: CH01-VOC-wyrt
wyrt (f.) & herb; plant & G sg wyrte\newline N pl wyrta & wyrt/tīd \\
% teaching: CH01-VOC-drygnes
drȳgnes (f.) & dry land & G sg drȳgnesse\newline N pl drȳgnessa & andlīcnes; S \\
% teaching: CH01-VOC-waeter
wæter (n.) & water & G sg wæteres\newline N pl wæteru & §4 \\
% teaching: CH01-VOC-andlicnes
andlīcnes (f.) & image & G sg andlīcnesse\newline N pl andlīcnessa & andlīcnes \\
% teaching: CH01-VOC-gelicnes
gelīcnes (f.) & likeness & G sg gelīcnesse\newline N pl gelīcnessa & andlīcnes \\
% teaching: CH01-VOC-leoht
lēoht (n.) & light & G sg lēohtes\newline N pl lēoht & hof/sǣd; S \\
% teaching: CH01-VOC-thiestru
ðīestru (P) & darkness & G sg —\newline N pl ðīestru & ðīestru; P; G pl ðīestra \\
% teaching: CH01-VOC-tacen
tācen (n.) & sign & G sg tācnes\newline N pl tācnu & §4 \\
% teaching: CH01-VOC-tid
tīd (f.) & season; time & G sg tīde\newline N pl tīda & wyrt/tīd \\
% teaching: CH01-VOC-gast
gāst (m.) & spirit & G sg gāstes\newline N pl gāstas & fisc \\
% teaching: CH01-VOC-lif
līf (n.) & life & G sg līfes\newline N pl līf & hof/sǣd; S \\
% teaching: CH01-VOC-gesceaft
gesceaft (f.) & creature; created being or thing & G sg gesceafte\newline N pl gesceafta & gesceaft \\
% teaching: CH01-VOC-mete
męte (m.) & food & G sg mętes\newline N pl mętas & męte; S \\
% teaching: CH01-VOC-gear
gēar (n.) & year & G sg gēares\newline N pl gēar & hof/sǣd \\
% teaching: CH01-VOC-mann
mann (m.) & man; human being & G sg mannes\newline N pl męnn & mann \\
% teaching: CH01-VOC-deor
dēor (n.) & animal; beast & G sg dēores\newline N pl dēor & hof/sǣd \\
% teaching: CH01-VOC-waestm
wæstm (m.) & fruit; produce & G sg wæstmes\newline N pl wæstmas & fisc \\
% teaching: CH01-VOC-faestnes
fæstnes (f.) & firmament & G sg fæstnesse\newline N pl fæstnessa & andlīcnes; S \\
% teaching: CH01-VOC-neowolnes
nēowolnes (f.) & abyss; the deep & G sg nēowolnesse\newline N pl nēowolnessa & andlīcnes; S \\
% teaching: CH01-VOC-sae
sǣ (m.) & sea & G sg sǣs\newline N pl sǣs & sǣ \\
% teaching: CH01-VOC-hwael
hwæl (m.) & whale & G sg hwæles\newline N pl hwalas & dæg/hwæl \\
% teaching: CH01-VOC-nieten
nīeten (n.) & animal; beast; in the plural, cattle & G sg nīetenes\newline N pl nīetenu & §4 \\
% teaching: CH01-VOC-anginn
anginn (n.) & beginning & G sg anginnes\newline N pl anginn & hof/sǣd \\
% teaching: CH01-VOC-god
God (m.) & God (proper divine name) & G sg Godes\newline N pl — & fisc; divine name, singular \\
% teaching: CH01-VOC-bradnes
brādnes (f.) & surface; face (of the deep) & G sg brādnesse\newline N pl brādnessa & andlīcnes; S \\
% teaching: CH01-VOC-gegaderung
gegaderung (f.) & gathering; collection; congregation (of waters) & G sg gegaderunge\newline N pl gegaderunga & gegaderung \\
% teaching: CH01-VOC-lyft
lyft (n.) & air & G sg lyftes\newline N pl lyft & hof/sǣd; S \\
% teaching: CH01-VOC-fraetwung
frætwung (f.) & array & G sg frætwunge\newline N pl frætwunga & gegaderung; S \\
% teaching: CH01-VOC-weorc
weorc (n.) & work & G sg weorces\newline N pl weorc & hof/sǣd \\
% teaching: CH01-VOC-cynn
cynn (n.) & kind & G sg cynnes\newline N pl cynn & hof/sǣd \\
% teaching: CH01-VOC-steorra
steorra (m.) & star & G sg steorran\newline N pl steorran & steorra/eorðe \\
% teaching: CH01-VOC-liehting
līehting (f.) & illumination & G sg līehtinge\newline N pl līehtinga & gegaderung; S \\
% teaching: CH01-VOC-hof
hof (n.) & dwelling & G sg hofes\newline N pl hofu & hof/sǣd \\
\end{longtable}
\par\endgroup

The other words required for composition are \textbf{and}, “and”; \textbf{on + dative}, “in/on” for location; \textbf{fram + dative}, “from” in separation; the article; and the indicative forms in §8. Verbs printed in frames and the additional help for the reading are supplied, not additional vocabulary to recall independently.

\section*{11. Forms and agreement}

The tasks in §§11–12 and §§14–17 are new practice. Give the forms selected in the tables; the key distinguishes valid source variants from the particular pattern requested.

\begin{samepage}
% item: CH01-F01
\itempara{1}{Decline fisc with its definite article in all four cases, singular and plural. Label every case and number.}

\end{samepage}

\begin{samepage}
% item: CH01-F02
\itempara{2}{Give nominative and accusative plurals of hof and sǣd. Explain the different endings.}

\end{samepage}

\begin{samepage}
% item: CH01-F03
\itempara{3}{Give genitive plural dæg and hwæl, dative singular mann, and genitive plural niht.}

\end{samepage}

\begin{samepage}
% item: CH01-F04
\itempara{4}{Give trēo genitive singular, dative singular and nominative plural; tācen genitive singular and nominative plural; wæter genitive singular and nominative plural; ǣfen and nīeten nominative plural.}

\end{samepage}

\begin{samepage}
% item: CH01-F05
\itempara{5}{Decline steorra in all four cases, singular and plural, then eorðe in all four singular cases.}

\end{samepage}

\begin{samepage}
% item: CH01-F06
\itempara{6}{Give accusative and genitive singular wyrt; accusative and genitive singular tīd; accusative singular gesceaft.}

\end{samepage}

\begin{samepage}
% item: CH01-F07
\itempara{7}{Give dative singular fæstnes, genitive plural andlīcnes, nominative plural gelīcnes, and dative plural gegaderung.}

\end{samepage}

\begin{samepage}
% item: CH01-F08
\itempara{8}{Locate ðȳ and ðon in the article table. Identify their case and singular genders; this is a reference-form identification.}

\end{samepage}

\subsection*{The copula}

Give the present and preterite indicative for each person and number. Use the English labels; no Old English personal pronoun is required.

\begin{samepage}
% item: CH01-COP01
\itempara{1}{First person singular.}

\end{samepage}

\begin{samepage}
% item: CH01-COP02
\itempara{2}{Second person singular.}

\end{samepage}

\begin{samepage}
% item: CH01-COP03
\itempara{3}{Third person singular.}

\end{samepage}

\begin{samepage}
% item: CH01-COP04
\itempara{4}{First person plural.}

\end{samepage}

\begin{samepage}
% item: CH01-COP05
\itempara{5}{Second person plural.}

\end{samepage}

\begin{samepage}
% item: CH01-COP06
\itempara{6}{Third person plural.}

\end{samepage}

\section*{12. Case in sentences}

Complete each frame with the noun and article phrases needed for the English sentence. The printed verb is supplied in its exact form; choose the noun cases, numbers, and articles. The prepositions in the frames use the constructions of §9.

\par\noindent\begin{minipage}{\linewidth}\raggedright
% item: CH01-CASE01
\itempara{1}{The man saw the bird.}

\continuation{\textbf{\_\_\_\_\_ geseah \_\_\_\_\_.}}

\continuation{Supplied: \textbf{geseah}, “saw,” third-person singular preterite.}

\end{minipage}\par

\par\noindent\begin{minipage}{\linewidth}\raggedright
% item: CH01-CASE02
\itempara{2}{The bird saw the man.}

\continuation{\textbf{\_\_\_\_\_ geseah \_\_\_\_\_.}}

\continuation{Supplied: \textbf{geseah}, “saw,” third-person singular preterite.}

\end{minipage}\par

\par\noindent\begin{minipage}{\linewidth}\raggedright
% item: CH01-CASE03
\itempara{3}{God created the heavens and the earth. Use masculine heofon for 'heavens'.}

\continuation{\textbf{\_\_\_\_\_ gescēop \_\_\_\_\_ and \_\_\_\_\_.}}

\continuation{Supplied: \textbf{gescēop}, “created,” third-person singular preterite.}

\end{minipage}\par

\par\noindent\begin{minipage}{\linewidth}\raggedright
% item: CH01-CASE04
\itempara{4}{The man saw the stars.}

\continuation{\textbf{\_\_\_\_\_ geseah \_\_\_\_\_.}}

\continuation{Supplied: \textbf{geseah}, “saw,” third-person singular preterite.}

\end{minipage}\par

\par\noindent\begin{minipage}{\linewidth}\raggedright
% item: CH01-CASE05
\itempara{5}{The man saw the fish in the sea.}

\continuation{\textbf{\_\_\_\_\_ geseah \_\_\_\_\_ on \_\_\_\_\_.}}

\continuation{Supplied: \textbf{geseah}, “saw,” third-person singular preterite. The fish is in the sea.}

\end{minipage}\par

\par\noindent\begin{minipage}{\linewidth}\raggedright
% item: CH01-CASE06
\itempara{6}{God separated the light from the darkness.}

\continuation{\textbf{\_\_\_\_\_ gedǣlde \_\_\_\_\_ fram \_\_\_\_\_.}}

\continuation{Supplied: \textbf{gedǣlde}, “separated/divided,” third-person singular preterite.}

\end{minipage}\par

\par\noindent\begin{minipage}{\linewidth}\raggedright
% item: CH01-CASE07
\itempara{7}{The bird saw the seed in the man's dwelling.}

\continuation{\textbf{\_\_\_\_\_ geseah \_\_\_\_\_ on \_\_\_\_\_.}}

\continuation{Supplied: \textbf{geseah}, “saw,” third-person singular preterite. The seed is in the dwelling.}

\end{minipage}\par

\par\noindent\begin{minipage}{\linewidth}\raggedright
% item: CH01-CASE08
\itempara{8}{The bird saw the man in the dwelling.}

\continuation{\textbf{\_\_\_\_\_ geseah \_\_\_\_\_ on \_\_\_\_\_.}}

\continuation{Supplied: \textbf{geseah}, “saw,” third-person singular preterite. The man is in the dwelling.}

\end{minipage}\par

\par\noindent\begin{minipage}{\linewidth}\raggedright
% item: CH01-CASE09
\itempara{9}{Complete the printed sealde frame to express: God gave food to the man. Identify the recipient and the thing given.}

\continuation{\textbf{\_\_\_\_\_ sealde \_\_\_\_\_ \_\_\_\_\_.}}

\continuation{Supplied: \textbf{sealde}, “gave,” third-person singular preterite. Put the recipient before the thing given in this frame.}

\end{minipage}\par

\par\noindent\begin{minipage}{\linewidth}\raggedright
% item: CH01-CASE10
\itempara{10}{Complete the printed sealde frame to express: The man gave seed to the bird. Identify the recipient and the thing given.}

\continuation{\textbf{\_\_\_\_\_ sealde \_\_\_\_\_ \_\_\_\_\_.}}

\continuation{Supplied: \textbf{sealde}, “gave,” third-person singular preterite. Put the recipient before the thing given in this frame.}

\end{minipage}\par

\section*{13. Cook's noun-phrase exercises}

Translate the following original prompts using the vocabulary of this unit. In Exercises 1–2, make the isolated phrases nominative and preserve the requested number. In Exercise 3, keep the head nominative; place its genitive dependent first, omitting the head's own article as explained in §9. The dependent retains its appropriate article. Use \textbf{dēor} for “animal/beast” and \textbf{nīeten} for plural “cattle”; \textbf{brādnes} means the surface in “face of the deep.”[14]

\subsection*{Group I, Exercise 1}

\begin{samepage}
% item: EX-I-01-01
\itempara{1}{The evening and the morning.}

\end{samepage}

\begin{samepage}
% item: EX-I-01-02
\itempara{2}{The day and the night.}

\end{samepage}

\begin{samepage}
% item: EX-I-01-03
\itempara{3}{The earth and the heaven.}

\end{samepage}

\begin{samepage}
% item: EX-I-01-04
\itempara{4}{The fish and the bird.}

\end{samepage}

\begin{samepage}
% item: EX-I-01-05
\itempara{5}{The grass and the seed.}

\end{samepage}

\begin{samepage}
% item: EX-I-01-06
\itempara{6}{The tree and the herb.}

\end{samepage}

\begin{samepage}
% item: EX-I-01-07
\itempara{7}{The dry land and the water.}

\end{samepage}

\begin{samepage}
% item: EX-I-01-08
\itempara{8}{The image and the likeness.}

\end{samepage}

\begin{samepage}
% item: EX-I-01-09
\itempara{9}{The light and the darkness (plur.).}

\end{samepage}

\begin{samepage}
% item: EX-I-01-10
\itempara{10}{The sign and the season.}

\end{samepage}

\begin{samepage}
% item: EX-I-01-11
\itempara{11}{The spirit and the life.}

\end{samepage}

\begin{samepage}
% item: EX-I-01-12
\itempara{12}{The creature and the food.}

\end{samepage}

\begin{samepage}
% item: EX-I-01-13
\itempara{13}{The year and the day.}

\end{samepage}

\begin{samepage}
% item: EX-I-01-14
\itempara{14}{The man and the animal.}

\end{samepage}

\begin{samepage}
% item: EX-I-01-15
\itempara{15}{The tree and the fruit.}

\end{samepage}

\begin{samepage}
% item: EX-I-01-16
\itempara{16}{The firmament and the abyss.}

\end{samepage}

\subsection*{Group I, Exercise 2}

\begin{samepage}
% item: EX-I-02-01
\itempara{1}{The days and the nights.}

\end{samepage}

\begin{samepage}
% item: EX-I-02-02
\itempara{2}{The fishes and the birds.}

\end{samepage}

\begin{samepage}
% item: EX-I-02-03
\itempara{3}{The trees and the herbs.}

\end{samepage}

\begin{samepage}
% item: EX-I-02-04
\itempara{4}{The images and the likenesses.}

\end{samepage}

\begin{samepage}
% item: EX-I-02-05
\itempara{5}{The signs and the seasons.}

\end{samepage}

\begin{samepage}
% item: EX-I-02-06
\itempara{6}{The years and the days.}

\end{samepage}

\begin{samepage}
% item: EX-I-02-07
\itempara{7}{The men and the animals.}

\end{samepage}

\begin{samepage}
% item: EX-I-02-08
\itempara{8}{The seas and the heavens.}

\end{samepage}

\begin{samepage}
% item: EX-I-02-09
\itempara{9}{The seeds and the fruits.}

\end{samepage}

\begin{samepage}
% item: EX-I-02-10
\itempara{10}{The evenings and the mornings.}

\end{samepage}

\begin{samepage}
% item: EX-I-02-11
\itempara{11}{The spirits and the men.}

\end{samepage}

\begin{samepage}
% item: EX-I-02-12
\itempara{12}{The whales and the cattle.}

\end{samepage}

\subsection*{Group I, Exercise 3}

\begin{samepage}
% item: EX-I-03-01
\itempara{1}{The beginning of the day.}

\end{samepage}

\begin{samepage}
% item: EX-I-03-02
\itempara{2}{The spirit of God.}

\end{samepage}

\begin{samepage}
% item: EX-I-03-03
\itempara{3}{The light of the life.}

\end{samepage}

\begin{samepage}
% item: EX-I-03-04
\itempara{4}{The life of the bird.}

\end{samepage}

\begin{samepage}
% item: EX-I-03-05
\itempara{5}{The food of the fish.}

\end{samepage}

\begin{samepage}
% item: EX-I-03-06
\itempara{6}{The firmament of the heaven.}

\end{samepage}

\begin{samepage}
% item: EX-I-03-07
\itempara{7}{The beasts of the earth.}

\end{samepage}

\begin{samepage}
% item: EX-I-03-08
\itempara{8}{The face of the deep.}

\end{samepage}

\begin{samepage}
% item: EX-I-03-09
\itempara{9}{The congregations of the waters.}

\end{samepage}

\begin{samepage}
% item: EX-I-03-10
\itempara{10}{The waters of the sea.}

\end{samepage}

\begin{samepage}
% item: EX-I-03-11
\itempara{11}{The fishes of the sea.}

\end{samepage}

\begin{samepage}
% item: EX-I-03-12
\itempara{12}{The birds of the air.}

\end{samepage}

\begin{samepage}
% item: EX-I-03-13
\itempara{13}{The array of the heavens and of the earth.}

\end{samepage}

\begin{samepage}
% item: EX-I-03-14
\itempara{14}{The works of God.}

\end{samepage}

\begin{samepage}
% item: EX-I-03-15
\itempara{15}{The kinds of the animals.}

\end{samepage}

\begin{samepage}
% item: EX-I-03-16
\itempara{16}{The stars of the heaven.}

\end{samepage}

\begin{samepage}
% item: EX-I-03-17
\itempara{17}{The seeds of the grass.}

\end{samepage}

\begin{samepage}
% item: EX-I-03-18
\itempara{18}{The illumination of the nights.}

\end{samepage}

\section*{14. Composition}

Use the selected forms and constructions without supplied verb frames.

\begin{samepage}
% item: CH01-COMP01
\itempara{1}{The bird is in the tree.}

\end{samepage}

\begin{samepage}
% item: CH01-COMP02
\itempara{2}{The birds were in the trees.}

\end{samepage}

\begin{samepage}
% item: CH01-COMP03
\itempara{3}{The seed is the bird's food.}

\end{samepage}

\begin{samepage}
% item: CH01-COMP04
\itempara{4}{The fish is an animal.}

\end{samepage}

\begin{samepage}
% item: CH01-COMP05
\itempara{5}{The signs were in the sky.}

\end{samepage}

\begin{samepage}
% item: CH01-COMP06
\itempara{6}{The herbs are in the man's dwelling.}

\end{samepage}

\begin{samepage}
% item: CH01-COMP07
\itempara{7}{The light of the night was in the sky.}

\end{samepage}

\begin{samepage}
% item: CH01-COMP08
\itempara{8}{The animals' food is the fruit of the trees.}

\end{samepage}

\section*{15. Short readings}

These two passages are new composition using the vocabulary and grammar of this unit.

\par\noindent\begin{minipage}{\linewidth}\raggedright
% item: CH01-MINI01
\itempara{1}{Read the complete passage without extra glosses. Translate it; identify whose dwelling is mentioned and explain the cases of mannes and hofe in its second sentence.}

\continuation{Se mann is on ðǣm hofe. Ðæt sǣd is on ðæs mannes hofe. Se fugol is on ðǣm trēowe. Ðæt sǣd is ðæs fugles męte.}

\end{minipage}\par

\par\noindent\begin{minipage}{\linewidth}\raggedright
% item: CH01-MINI02
\itempara{2}{Translate the passage. Rewrite it in the present. Explain the number of trēowu and the case of eorðan.}

\continuation{Ðā fiscas wǣron on ðǣm sǣ. Ðā fuglas wǣron on ðǣm trēowum. Ðā trēowu wǣron on ðǣre eorðan.}

\end{minipage}\par

\section*{16. The first day of Creation}

Cook, \emph{A First Book in Old English}, Reader I, pp.123–124. The following is a continuous extract through the close of the first day. Only note cues and printed line/page breaks have been removed.[15]

% source: FB-I-CREATION-FIRST-DAY-01
\begin{samepage}
\readingtext{On anginne gescēop God heofonan and eorðan. Sēo eorðe sōðlice wæs īdel and ǣmtig; and ðīestru wǣron ofer ðǣre nēowolnesse brādnesse; and Godes gāst wæs gefęred ofer wæteru. God cwæð ðā, “Geweorðe lēoht”; and lēoht wearð geworht. God geseah ðā ðæt hit gōd wæs; and hē gedǣlde ðæt lēoht fram ðǣm ðīestrum. And hēt ðæt lēoht Dæg, and ðā ðīestru Niht. Ðā wæs geworden ǣfen and morgen ān dæg.}
\end{samepage}

\subsection*{Supplied help}

% teaching: CH01-READ-SUPPORT CH01-GRAM-ORDER
Use this help for the questions and translation. It supplies the actual forms and constructions occurring here; their independent production is not required.

\par\begingroup
\fontsize{10.5}{13.2}\selectfont
\setlength{\LTleft}{0pt}\setlength{\LTright}{\fill}
\begin{longtable}{@{}>{\raggedright\arraybackslash}p{43mm}>{\raggedright\arraybackslash}p{79mm}@{}}
\toprule
\textbf{Expression} & \textbf{Local meaning and analysis} \\
\midrule
\endfirsthead
\toprule
\textbf{Expression} & \textbf{Local meaning and analysis} \\
\midrule
\endhead
\midrule
\multicolumn{2}{r}{\footnotesize Continued on the following page.}\\
\endfoot
\bottomrule
\endlastfoot
\textbf{On anginne} & “In the beginning”: a supplied temporal use of \textbf{on}, with dative \textbf{anginne}. \\
\textbf{gescēop} & “Created,” third-person singular preterite. In the opening, the subject \textbf{God} follows the verb; word order alone does not identify the object. \\
\textbf{heofonan} & “Heaven”: accusative \textbf{singular} of weak feminine \textbf{heofone}, a variant of the masculine \textbf{heofon} learned above. It is not masculine \textbf{heofon}'s plural. \\
\textbf{sōðlice; īdel and ǣmtig} & Connective “now/for its part”; “empty and void.” The two adjectives are predicate feminine nominative singular with \textbf{eorðe}. \\
\textbf{ofer ðǣre nēowolnesse brādnesse} & “Over the surface of the deep.” \textbf{Ofer} here governs dative \textbf{brādnesse}, the head. \textbf{Ðǣre nēowolnesse} is its genitive dependent; \textbf{ðǣre} belongs to that dependent. \\
\textbf{wæs gefęred ofer wæteru} & “Was borne/carried over the waters,” contextually “moved over the waters.” \textbf{Wæs gefęred} is a passive phrase; \textbf{gefęred} is a past participle. Here \textbf{ofer} has accusative plural \textbf{wæteru}, unlike the preceding dative construction. \\
\textbf{cwæð ðā} & “Then said”: \textbf{cwæð} is third-person singular preterite; this \textbf{ðā} is the adverb “then,” not the article. \\
\textbf{Geweorðe lēoht} & “Let light come to be.” \textbf{Geweorðe} is a finite third-person singular present subjunctive (also called optative), expressing a wish or command. \\
\textbf{lēoht wearð geworht} & “Light was made/came into being.” \textbf{Wearð} is a third-person singular preterite used here with past participle \textbf{geworht} in a passive phrase. \\
\textbf{geseah ðā ðæt hit gōd wæs} & “Then saw that it was good.” \textbf{Geseah} is third-person singular preterite; \textbf{ðā} is “then”; \textbf{ðæt} is the conjunction “that.” \textbf{Hit}, “it,” refers to neuter \textbf{lēoht}; \textbf{gōd} is a neuter nominative singular predicate adjective. The copula stands at the end of the embedded clause. \\
\textbf{hē gedǣlde} & “He separated”: \textbf{hē} is nominative singular, referring to God; \textbf{gedǣlde} is third-person singular preterite. The object and \textbf{fram} phrase use §9. \\
\textbf{hēt ðæt lēoht Dæg, and ðā ðīestru Niht} & “[He] called the light Day, and the darkness Night.” \textbf{Hēt} is third-person singular preterite, with God/he understood as subject. \textbf{Ðæt lēoht} and \textbf{ðā ðīestru} are accusative objects; \textbf{Dæg} and \textbf{Niht} are accusative names applied to them. Here \textbf{ðæt} and \textbf{ðā} are articles. \\
\textbf{Ðā wæs geworden ǣfen and morgen ān dæg.} & “Then there came to be evening and morning: one day.” \textbf{Wæs geworden} expresses a resulting state or completed coming-to-be, “had come to be”; \textbf{geworden} is a past participle. \textbf{Ān dæg} is “one day,” not the ordinal “first day.” Treat this closing formula as a whole. \\
\end{longtable}
\par\endgroup

\subsection*{Questions and translation}

\begin{samepage}
% item: CH01-READ01
\itempara{1}{Using the local support, identify the subject and two objects of the first sentence. Is heofonan plural? Translate the sentence.}

\end{samepage}

\begin{samepage}
% item: CH01-READ02
\itempara{2}{Describe the earth and the darkness. Explain wǣron, then identify the head and genitive dependent in the phrase after ofer.}

\end{samepage}

\begin{samepage}
% item: CH01-READ03
\itempara{3}{With the supplied passive phrase, state what God's spirit was doing and over what. Explain why wæteru cannot be the selected dative plural form.}

\end{samepage}

\begin{samepage}
% item: CH01-READ04
\itempara{4}{Does God's quoted speech report an existing fact or command/wish for something to come about? State what follows.}

\end{samepage}

\begin{samepage}
% item: CH01-READ05
\itempara{5}{In God geseah ðā ðæt hit gōd wæs, identify the references/functions of ðā, ðæt and hit using the support. Translate the clause.}

\end{samepage}

\begin{samepage}
% item: CH01-READ06
\itempara{6}{What did God separate, and from what? Identify the direct-object and separation phrases and their cases.}

\end{samepage}

\begin{samepage}
% item: CH01-READ07
\itempara{7}{Using the naming and closing-formula support, state the names given, the understood subject of hēt, and the number of days in the close.}

\end{samepage}

\begin{samepage}
% item: CH01-READ08
\itempara{8}{Translate the entire first-day passage from the printed Old English, using the supplied help. Preserve the sequence, command, naming and one-day close.}

\end{samepage}

\section*{17. Review}

\begin{samepage}
% item: CH01-REVIEW01
\itempara{1}{A learner writes *dægas and *hwælas for the plurals of dæg and hwæl, then changes wæstm to *wastmas. Correct the intended plurals and explain which vowel alternation applies.}

\end{samepage}

\begin{samepage}
% item: CH01-REVIEW02
\itempara{2}{In on ðæs mannes hofe, a learner takes ðæs to agree with hofe. Identify the head and the dependent, and give each case.}

\end{samepage}

\begin{samepage}
% item: CH01-REVIEW03
\itempara{3}{List the possible case/number values of eorðan in the selected paradigm. Which applies in on ðǣre eorðan, and which in God gescēop ðā eorðan ("God created the earth")? The supplied gescēop means 'created'.}

\end{samepage}

\begin{samepage}
% item: CH01-REVIEW04
\itempara{4}{Explain why the guided reading has ðīestru wǣron although ordinary English uses 'darkness was'. Then identify the cases in ðæt lēoht fram ðǣm ðīestrum.}

\end{samepage}

\clearpage
\section*{Source notes}

\begingroup\fontsize{10}{12.8}\selectfont

\itempara{1}{Cook, \emph{A First Book in Old English}, second edition (1895), §§2–4, pp.5–6 / PDF21–22. The orthography follows the selected edition; this is not a historical account of all spellings or a complete pronunciation lesson.}

\itempara{2}{Cook, §§38–42, pp.26–27 / PDF42–43. Tables composed from the checked paradigms are teaching models; they do not claim separate textual attestation of every form.}

\itempara{3}{Cook, §84, p.50 / PDF66. The masculine article is \textbf{se}; the source distinguishes emphatic \textbf{sē}. Instrumental syntax is reserved for later treatment.}

\itempara{4}{Cook, §§43–44, pp.28–30 / PDF44–46; especially §43.2 for \textbf{dæg/hwæl} and §43.4 for syncope and retained vowels.}

\itempara{5}{Cook, §47, pp.32–33 / PDF48–49. The tables retain the selected \textbf{trēowu, tācnu, wæteru, nīetenu, ǣfenu} patterns. The neuter \textbf{ǣfen} plural is derived from its checked pattern, not quoted from a newly located passage.}

\itempara{6}{Cook, §51, pp.34–36 / PDF50–52; §51.1 and the vocabulary's \textbf{wyrt} reference, p.324 / PDF340, support the endingless accusatives. The alternative N/A plural \textbf{-e} and singular \textbf{-a} of the \textbf{-ung} group are source variants.}

\itempara{7}{Cook, §53, p.37 / PDF53.}

\itempara{8}{Cook, §§46,52, pp.31–32,36 / PDF47–48,52; \textbf{mann/męnn}, §17, p.13 / PDF29; \textbf{niht}, Reader p.125 with note 28 / PDF141; \textbf{sǣ}, vocabulary p.324 / PDF340. Wright and Wright, \emph{Old English Grammar}, second edition (1914), §388, p.187 / PDF207, explicitly marks genitive plural \textbf{*sǣwa} as reconstructed. It is not tested here.}

\itempara{9}{Cook, §138, pp.78–79 / PDF94–95. The bounded indicative series is not the complete \textbf{wesan/bēon} system.}

\itempara{10}{Cook, §§149–150, p.88 / PDF104; §189, p.102 / PDF118.}

\itempara{11}{Cook, \emph{Exercises in Old English} (1895), Exercise 3 introductory note, p.2 / PDF10; \emph{First Book}, §153, pp.89–90 / PDF105–106, especially §153.i.}

\itempara{12}{Cook, \emph{First Book}, §§160,164(a), pp.93–94 / PDF109–110; §166, p.95 / PDF111; §§168,172, pp.96–97 / PDF112–113. \textbf{Sealde} is printed in §117, p.72 / PDF88. The other supplied frames occur in the assigned Creation extract. Examples in §9 are new composition; the separation sentence adapts the extract's subject \textbf{hē} to \textbf{God}.}

\itempara{13}{The noun choices follow Cook's English–Old English vocabulary and the checked \textbf{First Book} paradigms. \textbf{Wæstm}, with short \textbf{æ}, corrects Cook's \textbf{wǣstm}; this correction applies also to its inflected forms. The source spelling is preserved in the editorial record. \textbf{Gærs}, \textbf{męte}, and the selected noun genders are retained. \textbf{Dēor} is not restricted to modern English “deer”; \textbf{nīeten} supplies “cattle” in the plural.}

\itempara{14}{All 46 prompts are reproduced from Cook, \emph{Exercises}, Group I, Exercises 1–3, pp.1–2 / PDF9–10. Their English wording and original numbering are unchanged. Directions and answers are new teaching material.}

\itempara{15}{Cook, \emph{First Book}, Reader I, pp.123–124 / PDF139–140: the five prose lines on p.123 are editorially counted; the first three on p.124 are counted from the printed marginal line number 5. The extract preserves \textbf{heofonan, gefęred, Geweorðe}, and unlengthened \textbf{geseah}. Its apparatus and original lineation remain separately transcribed. All PDF page numbers in these notes are one-based and refer to the selected scans.}

\endgroup
\end{document}
